\chapter{Introdução}

\section{O que é Hayashi}

\hayashi{} é uma linguagem de programação interpretada projetada para análise
econométrica aplicada. O nome homenageia Fumio Hayashi, cujo livro \textit{Econometrics}
(2000) é referência canônica na área.

A linguagem tem três objetivos centrais:

\begin{enumerate}
  \item \textbf{Sintaxe familiar} para quem usa Stata — comandos como \lstinline{generate},
        \lstinline{summarize}, \lstinline{regress} funcionam como esperado.
  \item \textbf{Programação de propósito geral} — funções, closures, estruturas de controle,
        módulos. Não é apenas um conjunto de comandos estatísticos.
  \item \textbf{Desempenho sem configuração} — o interpretador é escrito em Rust; não há
        JVM, não há GC, não há dependências de runtime.
\end{enumerate}

\section{Filosofia de design}

\subsection*{Semântica clara antes de sintaxe concisa}

Quando uma operação muta um objeto, ela deve ser visivelmente imperativa. Quando retorna
um valor novo, deve parecer funcional. \hayashi{} distingue essas duas formas de maneira
consistente:

\begin{lstlisting}[caption={Forma imperativa × forma funcional}]
// imperativo: altera df no lugar (Stata-style)
generate df z = x^2

// funcional: retorna novo DataFrame, df inalterado
let df2 = generate(df, z = x^2)
\end{lstlisting}

\subsection*{Copy-on-Write para DataFrames}

DataFrames usam semântica de \textit{copy-on-write} (COW): atribuir um DataFrame a
outra variável cria um alias barato. A cópia real só acontece no momento da mutação.
Isso permite compartilhamento eficiente sem surpresas:

\begin{lstlisting}[caption={COW em ação}]
load "dados.csv" as df
let baseline = df      // alias, sem cópia

let transformado = df |> generate(z = x^2) |> filter(z > 10)
// baseline ainda não tem coluna z
\end{lstlisting}

\subsection*{Mensagens de erro como documentação}

Erros em \hayashi{} incluem o trecho de código que os causou, o número da linha e, quando
possível, uma sugestão de correção (\textit{did you mean?}). Um erro bem formatado é
mais útil que um manual.

\section{Hayashi e Stata}

\hayashi{} não é um clone de Stata. A sintaxe de comandos de dados é inspirada em Stata
porque essa sintaxe é boa — próxima da linguagem que economistas usam naturalmente.
Mas \hayashi{} é uma linguagem completa: tem funções de primeira classe, closures,
recursão, tratamento de erros estruturado.

A tabela abaixo resume as diferenças centrais:

\begin{center}
\begin{tabular}{lll}
\toprule
\textbf{Aspecto} & \textbf{Stata} & \textbf{Hayashi} \\
\midrule
Licença          & Proprietário    & GPL-3.0 \\
Funções          & Limitadas       & Primeira classe \\
Closures         & Não             & Sim \\
Tipos            & Implícitos      & Explícitos em funções \\
DataFrames       & Um por vez      & Múltiplos, COW \\
Modularização    & \texttt{do}     & \lstinline{source} \\
Desempenho       & C              & Rust \\
\bottomrule
\end{tabular}
\end{center}

\section{Hayashi e R / Python}

R e Python têm ecossistemas imensos. \hayashi{} não compete com eles no volume de
pacotes disponíveis — compete na clareza da sintaxe para a tarefa específica de
econometria aplicada. Um script de regressão com dados de painel em \hayashi{} tem
menos ruído visual que o equivalente em qualquer outra linguagem.

\section{Programa de validação empírica}

\hayashi{} inclui um programa de validação reprodutível no diretório
\texttt{validation/}. Cada caso executa um script \hayashi{} e o compara,
 dentro de tolerâncias declaradas, com implementações de referência em R e
Python (statsmodels). Os casos usam tanto datasets reais públicos quanto DGPs
simulados retirados dos capítulos deste livro.

Para executar a validação completa:

\begin{lstlisting}[caption={Executar o programa de validação}]
python -m venv validation/.venv
validation/.venv/bin/pip install -r validation/requirements.txt
hay validate
\end{lstlisting}

O estado atual de cada caso está registrado em \texttt{validation/MATRIX.md}.

\section{Estrutura deste livro}

Este volume cobre a Parte~I: a linguagem em si.

\begin{description}
  \item[Capítulos 2--3] Instalação e uso do REPL.
  \item[Capítulos 4--12] Sistema de tipos, variáveis, operadores, controle de fluxo,
        funções, erros, strings, listas e dicionários.
  \item[Capítulos 13--15] DataFrames, manipulação de dados e o operador pipe.
  \item[Capítulo 16] Modularização com \lstinline{source}.
\end{description}

A Parte~II cobrirá estatística descritiva, estimadores de OLS, IV, painel, séries de
tempo, modelos de variável dependente limitada e testes de hipótese.
